\begin{textsample}{POS Dim 3 – mt – Score 24.00 – mt/mt\_inf\_15.txt}  \label{ex:f3_pos_078}
2.2 Corpo, \textbf{Gestos} e Movimentos ( CG ) A \textbf{criança} na sua integralidade, conhece e explora o mundo por meio da linguagem corporal, sendo manifestada mediante aos \textbf{gestos}, expressão facial, \textbf{mímicas}, deslocamentos espaciais, manipulação e exploração dos objetos, das \textbf{brincadeiras} e da sua cultura, expressando suas \textbf{vontades} e emoções, vivenciando diferentes experiências.
Ao longo da história, a cultura escolar foi construída com a ideia de que o corpo é um elemento que atrapalha a aprendizagem, alegando que as \textbf{crianças}, em especial as da Educação \textbf{Infantil}, com argumentos que as \textbf{crianças} “ movimentam -se demais, não ficam quietas para desenhar, pintar, \textbf{conversar} ou \textbf{ouvir} histórias ”.
Nessa concepção o corpo tornou -se o recurso usado pelos educadores para contenção, higienização e disciplinamento dos movimentos \textbf{infantis}, sendo explorado apenas para o ensino de técnicas, enfatizando o segurar o \textbf{lápis} ou os talheres de forma correta, sentar e permanecer imóvel para fazer uma atividade, ou ainda relaxar para a realização de tarefas consideradas “ mais nobres ” : \textbf{cópias} de letras, memorização de palavras e outras.
Pesquisas apontam que rígidas restrições impostas ao movimento das \textbf{crianças} criam uma atmosfera de tensão e conflito entre elas e o professor, incompatível com um pleno desenvolvimento tanto da \textbf{criança} quanto do professor.
Na contemporaneidade, faz -se necessário conceber o corpo das \textbf{crianças} como um integrante privilegiado das práticas pedagógicas orientadas para a interação e criação com \textbf{parceiros} na Educação \textbf{Infantil}.
Na primeira \textbf{infância}, o corpo é o instrumento expressivo e comunicativo por excelência, em que as \textbf{crianças} por meio de \textbf{gestos}, expressões faciais e movimentos corporais, desde o \textbf{nascimento} e ao longo da vida humana, exploram o ambiente, expressam seus sentimentos e \textbf{vontades}, interagem e se comunicam com seus \textbf{parceiros}.
Na etapa da Educação \textbf{Infantil} é necessário valorizar o corpo e o movimento, relacionando -os aos objetivos e ações pedagógicas e promovendo o desenvolvimento dos aspectos que envolvam as : - Habilidades motoras – força, \textbf{equilíbrio}, flexibilidade, coordenação fina e ampla, lateralidade, entre outras ; - Habilidades comportamentais – desinibição, socialização, conceito de saúde ( qualidade de vida, tempo livre, lazer ), vivências emocionais etc. ; - Habilidades expressivas – fluência verbal, ritmo, expressão \textbf{dramática}, dicção e destreza manual.
Compreende -se que na Educação \textbf{Infantil}, o corpo é o ponto de partida para o desenvolvimento do trabalho pedagógico, visto que por meio da estimulação e desafios corporais, a \textbf{criança} irá apropriar -se dos seus sentidos e suas funções, manifestando suas preferências ou recusas.
Ao professor cabe assegurar em suas ações pedagógicas o \textbf{cuidado} físico, o desenvolvimento motor, a ampliação de repertório de \textbf{gestos}, o uso do corpo em diferentes espaços, promovendo a emancipação e a liberdade, evitando seu cerceamento em situações individuais e coletivas.
O Estado de Mato Grosso concebe o movimento, o \textbf{gesto} e o corpo como manifestações de linguagem, à medida que lhes são possibilitadas a progressão de suas competências corporais, por meio dos direitos de aprendizagens de conviver, \textbf{brincar}, participar, expressar, explorar e conhecer -se, agindo no ambiente pelo movimento, conhecendo o próprio corpo, expressando -se e interagindo por meio de jogos, \textbf{brincadeiras}, danças e dramatizações entre outros.
Na BNCC ( 2017 ) a ementa desse campo se apresenta da seguinte maneira : > Com o corpo ( por meio dos sentidos, \textbf{gestos}, movimentos impulsivos ou intencionais, coordenados ou espontâneos ), as \textbf{crianças}, desde cedo, exploram o mundo, o espaço e os objetos do seu entorno, estabelecem relações, expressam -se, \textbf{brincam} e produzem conhecimentos sobre si, sobre o outro, sobre o universo social e cultural, tornando -se, progressivamente, conscientes dessa corporeidade.
Por meio das diferentes linguagens, como a música, a dança, o teatro, as \textbf{brincadeiras} de faz de conta, elas se comunicam e se expressam no entrelaçamento entre corpo, emoção e linguagem.
As \textbf{crianças} conhecem e reconhecem as sensações e funções de seu corpo e, com seus \textbf{gestos} e movimentos, identificam suas potencialidades e seus limites, desenvolvendo, ao mesmo tempo, a consciência sobre o que é seguro e o que pode ser um risco à sua integridade física.
Na Educação \textbf{Infantil}, o corpo das \textbf{crianças} ganha \textbf{centralidade}, pois ele é o \textbf{partícipe} privilegiado das práticas pedagógicas de \textbf{cuidado} físico, orientadas para a emancipação e a liberdade, e não para a submissão.
Assim, a \textbf{instituição} escolar precisa promover oportunidades ricas para que as \textbf{crianças} possam, sempre animadas pelo espírito \textbf{lúdico} e na interação com seus \textbf{pares}, explorar e vivenciar um amplo repertório de movimentos, \textbf{gestos}, olhares, \textbf{sons} e \textbf{mímicas} com o corpo, para descobrir variados modos de ocupação e uso do espaço com o corpo ( tais como sentar com apoio, rastejar, engatinhar, escorregar, caminhar \textbf{apoiando} -se em berços, mesas e cordas, saltar, escalar, equilibrar -se, correr, dar cambalhotas, alongar -se etc. ) ( BRASIL, 2017, p. 38-39 ).

% matched lemmas: apoiar, brincadeira, brincar, centralidade, conversar, criança, cuidado, cópia, dramático, equilíbrio, gesto, infantil, infância, instituição, lápis, lúdico, mímico, nascimento, ouvir, par, parceiro, partícipe, som, vontade
\end{textsample}
